\title{FlashAttention-3: Fast and Accurate Attention with Asynchrony and Low-precision}

\begin{abstract}
The attention mechanism is fundamental to the Transformer architecture, yet it remains a primary performance constraint in large-scale language models and contexts that demand lengthy inputs. \textsc{FlashAttention} previously accelerated attention computation on GPUs by reducing memory traffic. Despite these advances, recent hardware innovations have not been fully harnessed; for instance, \textsc{FlashAttention-2} manages only 35\% efficiency on NVIDIA’s H100 GPU. This work introduces three techniques tailored for the Hopper GPU series: we utilize the asynchronous features of Tensor Cores and TMA to (1) achieve concurrency between computation and data transfer through warp specialization, (2) intertwine block-level matrix multiplications with softmax processing, and (3) employ block quantization alongside incoherent execution, enabled by native support for FP8 precision. Our approach, \textsc{FlashAttention-3}, demonstrates a significant speedup on H100 GPUs by a factor of 1.5-2.0$\times$, delivering FP16 performance up to 740 TFLOPs/s (amounting to 75\% hardware utilization) and FP8 rates approaching 1.2 PFLOPs/s. Empirically, FP8 \textsc{FlashAttention-3} produces 2.6$\times$ smaller numerical errors relative to standard FP8 attention baselines.
\end{abstract}

\section{Introduction}
\label{sec:intro}

Within the Transformer framework~\citep{vaswani2017attention}, the attention mechanism is recognized as the dominant computational constraint, chiefly because calculating the self-attention between queries and keys exhibits a quadratic time complexity with respect to sequence length. Extending attention mechanisms to accommodate longer contexts presents opportunities for enhancing model capabilities, such as enabling comprehension and inference across multiple extended documents~\citep{guo2021longt5,shaham2022scrolls,peng2023yarn} and large-scale code repositories~\citep{roziere2023code, li2023starcoder}. Additionally, it opens pathways to process new data modalities, including high-resolution imagery~\citep{chen2022scaling}, audio streams~\citep{gulati2020conformer}, and video~\citep{ho2022video}, as well as to address novel use-cases, such as modeling user behavior across extensive interaction histories~\citep{sun2019bert4rec} or managing agent tasks over extended timelines~\citep{yao2022react}. As a result, there has been considerable focus on accelerating attention under long-context conditions, utilizing methods such as approximate algorithms~\citep{katharopoulos2020transformers,choromanski2020rethinking, tay2020efficient}, advances in software optimization (\citep{rabe2021self,dao2022flashattention,kwon2023efficient}), and the development of alternative model architectures~\citep{peng2023rwkv,sun2023retentive,gu2023mamba}.

This study extends the foundational contributions by \citet{dao2022flashattention}, who devised exact-attention algorithms that thoughtfully account for the GPU’s hardware and execution model at the design phase. The seminal work \citep{dao2022flashattention} introduced \textsc{FlashAttention}, leveraging a novel tiling approach to efficiently parallelize attention computations, and streamlining the process by consolidating all attention steps into a single GPU kernel, thereby avoiding inefficient intermediate global memory accesses. Subsequently, \citet{dao2023flashattention2} revised this algorithm into \textsc{FlashAttention-2}, further enhancing parallelism across the sequence length axis and reconfiguring the forward pass to operate over blocks of key and value matrices—improving workload allocation and resource utilization on the GPU. Despite these improvements, we note that \textsc{FlashAttention-2} demonstrates suboptimal efficiency on modern GPU architectures, especially compared to high-performance matrix-multiplication (GEMM) routines, yielding, for example, only 35\% utilization as opposed to the 80-90\% achieved on the Hopper H100 GPU. This discrepancy can be partly attributed to implementation details, such as relying on Ampere rather than Hopper-tailored instructions for Tensor Cores. Prior works like ThunkerKitten~\citep{spector2024thunder} and cuDNN 9~\citep{cudnn9} have demonstrated that incorporating Hopper-specific instructions and tile-oriented designs can accelerate attention mechanisms and streamline their coding.

At a deeper level, the algorithm underpinning \textsc{FlashAttention-2} operates within a straightforward synchronous framework and does not incorporate asynchronous execution or leverage low-precision computation in its architecture. Asynchronous operation originates from specialized hardware elements that optimize key tasks in machine learning pipelines: dedicated modules for matrix multiplication (Tensor Cores) and memory access (Tensor Memory Accelerator -- TMA), which function independently of general CUDA cores responsible for logical, integer, and floating-point arithmetic. The adoption of reduced precision formats, like FP8 in Hopper and FP4 in Blackwell—extending the legacy of FP16 (introduced with Pascal in 2017) and BF16 (featured in Ampere since 2020)—has demonstrated the ability to substantially increase computational throughput, effectively doubling or quadrupling performance while maintaining constant power consumption and die size. In \cref{subsec:hardware}, we survey the functionalities presented by Hopper along these axes. 
The core engineering problem is to rearchitect \textsc{FlashAttention-2} so that it utilizes these advanced hardware traits: asynchronous processing demands concurrent execution of matmul and softmax operations, despite dependency constraints between their outputs, while integrating low-precision arithmetic necessitates strategies to mitigate quantization artifacts, particularly for rare features in LLMs~\citep{dettmers2208llm, sun2024massive}.

For this purpose, we introduce \textsc{FlashAttention-3}, presenting and integrating three novel concepts to enhance efficiency on the latest GPU architectures:\footnote{Our experiments are focused on NVIDIA's Hopper architecture, yet the underlying algorithm is applicable to any GPU design that offers strong support for asynchronous computation and low-precision operations.}

\begin{enumerate}

\item \textbf{Producer-Consumer asynchrony:} We propose a warp-specialized software pipelining strategy that leverages asynchronous data transfers and Tensor Core computations by assigning distinct warps as either producers or consumers of data. This division enables the algorithm to further conceal latencies associated with memory accesses and instruction issuance.

\item \textbf{Hiding softmax under asynchronous block-wise GEMMs:} We concurrently execute the non-GEMM operations in softmax, such as floating-point multiply-add and exponentiation, alongside asynchronous WGMMA instructions for GEMM. To facilitate this overlap, we restructure the \textsc{FlashAttention-2} algorithm so as to eliminate certain sequential bottlenecks between softmax and GEMM phases. For instance, in our two-stage variant, softmax computation occurs on one block of the score matrix at the same time that WGMMA asynchronously computes the succeeding block.

\item \textbf{Hardware-accelerated low-precision GEMM:} We modify the forward pass algorithm to make use of FP8 Tensor Cores for GEMM operations, which results in almost a twofold increase in observed TFLOPs/s. Accommodating this requires handling distinct memory layouts for FP32 accumulators and FP8 operands as prescribed by WGMMA. We employ block quantization techniques and incoherent processing in order to minimize accuracy degradation inherent to FP8 precision.
\end{enumerate}

For empirical assessment, we conduct comprehensive benchmarks of \textsc{FlashAttention-3} on the H100 SXM5 GPU, spanning various parameter configurations. Our results reveal the following: (1) FP16 demonstrates a 1.5–2.0$\times$ performance improvement in the forward pass over \textsc{FlashAttention-2} (with throughput peaking at 740 TFLOPs/s), and a 1.5–1.75$\times$ gain during the backward pass; (2) FP8 attains nearly 1.2 PFLOPs/s; (3) At extended sequence lengths, FP16 exhibits superior performance, and FP8 shows comparable results\footnote{Specifically, FP8 in \textsc{FlashAttention-3} is optimal for head dimension 64, while for head dimensions of 128 and 256 it matches performance in non-causal masking scenarios, and lags somewhat with causal masking.} when measured against NVIDIA's cutting-edge cuDNN attention implementation. Additionally, FP16 in \textsc{FlashAttention-3} preserves the same level of numerical error as \textsc{FlashAttention-2} and outperforms the standard attention method, owing to FP32 storage of intermediate computations like softmax scaling. Further, FP8 in \textsc{FlashAttention-3}—utilizing block quantization and incoherent computation—provides 2.6$\times$ greater accuracy relative to standard attention approaches employing per-tensor quantization, particularly when handling outlier features.

\textsc{FlashAttention-3} has been released as open-source under a flexible license\footnote{\textsc{FlashAttention-3} can be accessed at \texttt{https://github.com/Dao-AILab/flash-attention}}. We aim to incorporate it into PyTorch and Hugging Face libraries to maximize accessibility for researchers and developers.

\section{Background: Multi-Head Attention and GPU Characteristics}
\label{sec:background}

\subsection{Multi-Head Attention}
\label{subsec:multi_head_attn}

Let $\mathbf{Q}, \mathbf{K}, \mathbf{V} \in \mathbb{R}^{N \times d}$ be the query, key and value input sequences associated to a single head, where $N$ is the sequence length and $d$ is the head dimension. Then the attention output $\mathbf{O}$ is computed as:

\begin{equation*}
  \mathbf{S} = \alpha \mathbf{Q} \mathbf{K}^\top \in \mathbb{R}^{N \times N}, \quad \mathbf{P} = \mathrm{softmax}(\mathbf{S}) \in \mathbb{R}^{N \times N}, \quad \mathbf{O} = \mathbf{P}\mathbf{V} \in \mathbb{R}^{N \times d},
\end{equation*}

Here, the $\mathrm{softmax}$ function operates across each row, and conventionally, the scaling parameter is chosen as $\alpha = 1/\sqrt{d}$. To mitigate potential numerical instability arising from exponentiation, it is standard practice to subtract $\mathrm{rowmax}(\mathbf{S})$ from $\mathbf{S}$. In the context of multi-head attention (MHA), individual heads utilize distinct projection matrices for queries, keys, and values. The entire computation can be executed in parallel across all heads and batches, resulting in the complete output tensor.

Now let $\phi$ be a scalar loss function and let $\mathbf{d}(-) = \partial \phi / \partial (-)$ be notation for the gradient. Given the output gradient $\mathbf{dO} \in \mathbb{R}^{N \times d}$, we compute $\mathbf{dQ}$, $\mathbf{dK}$, and $\mathbf{dV}$ according to the chain rule as follows:

\begin{align*}
  \mathbf{dV} &= \mathbf{P}^\top \mathbf{dO} \in \mathbb{R}^{N \times d} \\
  \mathbf{dP} &= \mathbf{dO} \mathbf{V}^\top \in \mathbb{R}^{N \times N} \\
  \mathbf{dS} &= \mathrm{dsoftmax} (\mathbf{dP}) \in \mathbb{R}^{N \times N} \\
  \mathbf{dQ} &= \alpha \mathbf{dS} \mathbf{K} \in \mathbb{R}^{N \times d} \\
  \mathbf{dK} &= \alpha \mathbf{dS}^\top \mathbf{Q} \in \mathbb{R}^{N \times d},
\end{align*}

In this context, $\mathbf{d}s = (\mathrm{diag}(p) - p p^\top)\mathbf{d}p$ holds when $p = \mathrm{softmax}(s)$, considering $p$ as dependent on the vector $s$. We denote the row-wise application of this formula as $\mathrm{dsoftmax}(\mathbf{dP})$. For the backward pass of MHA, this operation can also be executed in parallel over all heads and batches.

\subsection{GPU hardware characteristics and execution model}
\label{subsec:hardware}

This section outlines components of the GPU execution model pertinent to \textsc{FlashAttention-3}, emphasizing how these manifest within the NVIDIA Hopper architecture as a specific example.

\paragraph{Memory hierarchy:} The GPU architecture structures its memory as a series of data localities, where increased capacity generally corresponds to decreased bandwidth (\cref{tab:gpu-hierarchy})\footnote{\citet{luo2024benchmarking} documents a shared memory bandwidth of 128 bytes per clock cycle per SM. By considering all 132 SMs operating at the boost frequency of 1830 MHz, this value is accordingly scaled.}. The global memory (GMEM)—alternatively referred to as HBM—consists of off-chip DRAM that is universally accessible by all streaming multiprocessors (SMs). When data resides in GMEM, it is automatically brought into a unified on-chip L2 cache for faster access. At the next level, each SM is equipped with its own limited-size, explicitly managed, and highly banked cache known as shared memory (SMEM). The hierarchy culminates with the register file present within each SM.

\paragraph{Thread hierarchy:} In the GPU programming paradigm, execution entities known as threads are structured into nested logical groups. Starting at the lowest granularity, threads are organized into warps, each containing 32 threads, which are further aggregated into warpgroups composed of 4 consecutive warps. Above these, threads are arranged in threadblocks (also referred to as cooperative thread arrays or CTAs), which can be combined into threadblock clusters (specific to Hopper architectures), and ultimately into grids at the highest level.

The two hierarchies exhibit a strong interconnection. Threads belonging to the same CTA are assigned simultaneously to the same SM, while CTAs within an identical cluster are scheduled together on a single GPC. All threads inside a CTA have direct access to SMEM, in contrast to RMEM, which provides each thread with up to 256 private registers.

\paragraph{Asynchrony and warp-specialization:}

GPUs function as throughput-oriented processors, utilizing parallelism and asynchronous operations to mitigate delays associated with memory access and computation. Hopper introduces the Tensor Memory Accelerator (TMA) as a specialized hardware mechanism for asynchronous data transfers between GMEM and SMEM \cite[\S7.29]{cuda}. Additionally, in contrast to earlier architectures like Ampere, Hopper’s Tensor Core—accessible through the warpgroup-level WGMMA instruction \cite[\S9.7.14]{ptx}—operates asynchronously, enabling it to fetch input data directly from shared memory.

Enabling hardware-level asynchrony facilitates the creation of warp-specialized kernels, in which warps within a CTA are assigned exclusively as either producers or consumers, and thus are dedicated solely to either data transfer or computational tasks. This specialization generally enhances the compiler's capability to produce efficient instruction sequences \citep{warp-specialization-2011}. Furthermore, Hopper provides support for flexible register allocation among warpgroups through the \verb|setmaxnreg| instruction \cite[\S9.7.17.1]{ptx}, enabling warps involved in MMAs to utilize a greater portion of RMEM compared to those solely responsible for TMA operations (which require only a single thread).

\paragraph{Low-precision number formats:}
\label{sec:low-precision-gpu}
Contemporary GPUs incorporate dedicated hardware designed to boost performance for low-precision arithmetic. As an illustration, the WGMMA instruction is capable of utilizing the FP8 Tensor Cores on the Hopper architecture, achieving double the throughput per SM relative to computations carried out with FP16 or BF16.

To effectively utilize FP8 WGMMA, it is necessary to be familiar with the operand layout restrictions. In the context of a GEMM operation $A \times B^{\top}$ where $A$ is an $M\times K$ matrix and $B$ is an $N\times K$ matrix, an operand is referred to as \emph{mn-major} if elements are stored contiguously along the outer $M$ or $N$ dimension, while an operand is described as \emph{k-major} when contiguity occurs in the inner $K$ dimension. For FP16 WGMMA, operands residing in SMEM may be provided in either mn-major or k-major formats. In contrast, FP8 WGMMA restricts inputs solely to the k-major configuration. Additionally, in scenarios like attention where consecutive GEMM operations need to be fused in a single kernel, differing layouts between FP32 accumulators and FP8 operands complicate the execution of dependent FP8 WGMMAs.

Within the realm of attention mechanisms, these constraints related to layout necessitate adjustments to the construction of an FP8 algorithm, as detailed in \cref{sec:algofp8}.

\subsection{Standard Attention and Flash Attention} In line with \citet{dao2022flashattention}, we use the term \textbf{standard attention} to refer to a GPU-based attention implementation where the intermediate matrices $\mathbf{S}$ and $\mathbf{P}$ are explicitly stored in high-bandwidth memory (HBM). The principal innovation behind \textsc{FlashAttention} involves exploiting a local softmax reduction to circumvent the costly memory operations typically required for these intermediate products, instead unifying the attention computation into a single kernel. This notion of a local softmax is operationalized by lines \ref{code-ws:softmax_start}-\ref{code-ws:softmax_end} within the consumer mainloop described in \cref{alg:flash3_wgmma_ws_only}, including the scaling applied to partitions of $\mathbf{O}$. A straightforward derivation demonstrating that this strategy accurately yields $\mathbf{O}$ is presented in \cite[\S 2.3.1]{dao2023flashattention2}.

\section{FlashAttention-3: Algorithm}
\label{sec:algo}

This section outlines the workings of the \textsc{FlashAttention-3} algorithm. The discussion prioritizes the forward pass, while the backward pass is detailed separately in~\cref{sec:algo_ws_bwd}. We begin by illustrating the incorporation of warp-specialization and a circular SMEM buffer within the foundational framework of \textsc{FlashAttention-2}. Next, we detail how leveraging the asynchronicity inherent to WGMMA enables construction of a 2-stage pipeline that overlaps GEMM and softmax computations. Lastly, we cover the adjustments necessary for FP8, addressing both layout alignment and improvements to accuracy through block quantization and incoherent operations.

\subsection{Producer-Consumer asynchrony through warp-specialization and pingpong scheduling}
\label{sec:algo_ws}

\paragraph{Warp-specialization}
Similar to \textsc{FlashAttention-2}, the forward pass in \textsc{FlashAttention-3} is highly parallelizable across batch size, head count, and query sequence length. Consequently, a CTA-level perspective of the procedure is appropriate, wherein the algorithm processes a tile $\mathbf{Q}_i$ of the query matrix to derive the associated tile $\mathbf{O}_i$ of the output. To streamline exposition, we initially describe the warp-specialization approach employing a circular SMEM buffer, excluding the additional GEMM-softmax overlap mechanism. Define $d$ as the head dimension, $N$ as the sequence length, and select a query block size $B_r$, thereby partitioning $\mathbf{Q}$ into $T_r = \lceil \frac{N}{B_r} \rceil$ blocks, namely $\mathbf{Q}_1, .., \mathbf{Q}_{T_r}$.

\begin{algorithm}[H]
    \caption{\small\label{alg:flash3_wgmma_ws_only}\textsc{FlashAttention-3} forward pass \textbf{excluding} intra-consumer overlapping -- CTA perspective}
    \begin{algorithmic}[1]
\REQUIRE Matrices $\mathbf{Q}_i \in \mathbb{R}^{B_r \times d}$ and $\mathbf{K}, \mathbf{V} \in \mathbb{R}^{N \times d}$ located in HBM, key block size $B_c$, with total block count $T_c = \lceil \frac{N}{B_c} \rceil$.
\STATE Instantiate pipeline object to coordinate barrier synchronization via a circular SMEM buffer with $s$ stages.
\IF {executing in producer warpgroup}
\STATE Free a designated number of registers.
\STATE Initiate loading $\mathbf{Q}_i$ from HBM into shared memory.
\STATE When load completes, notify consumers regarding $\mathbf{Q}_i$ availability.
\FOR{$0 \le j < T_c$}
    \STATE Await consumption of the $(j\,\%\,s)$th buffer stage.
    \STATE Load $\mathbf{K}_j$ and $\mathbf{V}_j$ from HBM into shared memory at stage $(j\,\%\,s)$.
    \STATE Following load completion, signal to consumers that $\mathbf{K}_j$, $\mathbf{V}_j$ are ready.
\ENDFOR
\ELSE \STATE Allocate a specified number of registers, determined by the consumer warp count.
\STATE Locally, set $\mathbf{O}_i = (0) \in \mathbb{R}^{B_r \times d}$; initialize $\ell_i = (0)$, $m_i = (-\infty)$, both in $\mathbb{R}^{B_r}$.
\STATE Wait for completion of the $\mathbf{Q}_i$ transfer to shared memory.
\FOR{$0 \le j < T_c$}
\STATE Wait for $\mathbf{K}_j$ transfer to shared memory to finish.
\STATE Calculate $\mathbf{S}_i^{(j)} = \mathbf{Q}_i \mathbf{K}_j^T$ (SS-GEMM). Commit computation and wait. \label{alg:ws_only_gemm1}
\STATE Set $m_i^{\mathrm{old}} = m_i$, update $m_i = \mathrm{max}(m_i^{\mathrm{old}}, \mathrm{rowmax}(\mathbf{S}_i^{(j)}))$. \label{code-ws:softmax_start}
\STATE Evaluate $\widetilde{\mathbf{P}}_i^{(j)} = \mathrm{exp}(\mathbf{S}_i^{(j)} - m_i)$, update $\ell_i = \mathrm{exp}(m_i^{\mathrm{old}} - m_i) \ell_i + \mathrm{rowsum}(\widetilde{\mathbf{P}}_i^{(j)})$. \label{code-ws:softmax_end}
\STATE Wait for $\mathbf{V}_j$ to finish loading into shared memory.
\STATE Update $\mathbf{O}_i = \mathrm{diag}(\mathrm{exp}(m_i^{\mathrm{old}} - m_i))^{-1} \mathbf{O}_i + \widetilde{\mathbf{P}}_i^{(j)} \mathbf{V}_j$ (RS-GEMM). Commit and synchronize. \label{alg:ws_only_gemm2}
\STATE Release buffer stage $(j\,\%\,s)$ for use by the producer.
\ENDFOR
\STATE Finalize $\mathbf{O}_i = \mathrm{diag}(\ell_i)^{-1} \mathbf{O}_i$ and set $L_i = m_i + \log(\ell_i)$.
\STATE Store $\mathbf{O}_i$ and $L_i$ into HBM, writing as the $i$th blocks of $\mathbf{O}$ and $L$.
\ENDIF
\end{algorithmic}
\end{algorithm}

In our deployment of \cref{alg:flash3_wgmma_ws_only} on Hopper, memory (de)allocation tasks are handled with \verb|setmaxnreg|, while TMA is used to load $\mathbf{Q}_i$ and the collection $\{ \mathbf{K}_j, \mathbf{V}_j \}_{0 \leq j < T_c}$. GEMM operations within the consumer mainloop are performed via WGMMA, utilizing the SS or RS prefix to specify whether the initial operand resides in shared memory or the register file, respectively. Understanding the execution sequence of \cref{alg:flash3_wgmma_ws_only}, it is important to recognize that TMA loads proceed independently, with no blocking on completion of prior loads due to their asynchronous nature. Additionally, throughout the first $s$ passes of the producer mainloop—while the buffer is being populated—no waits are invoked.

\paragraph{Pingpong scheduling} The interplay between the asynchronous execution of WGMMA and TMA and warp specialization enables the possibility of concurrently computing the softmax for one warpgroup while another warpgroup handles GEMM. This parallelism is motivated by the significant disparity in throughput between matmul and other operations on current hardware accelerators. For instance, the H100 SXM5 GPU delivers 989 TFLOPS for FP16 matmul, yet only achieves 3.9 TFLOPS for special function operations like exponential\footnote{According to the CUDA programming guide, each streaming multiprocessor (SM) is capable of executing 16 special function operations per clock cycle. To estimate throughput, this value is multiplied by 132 SMs and the 1830 MHz clock frequency, resulting in 3.9 TFLOPS for special functions.}, which are essential for computing softmax. In an FP16 attention forward pass with a head size of 128, matmul computations entail 512 times more FLOPS than exponential operations; however, exponentials are performed with 256 times lower throughput, implying that the exponential computations can occupy up to half of the processing cycle compared to matmul. This imbalance becomes even more pronounced for FP8, as matmul throughput is doubled while the throughput for exponential operations remains unchanged.

The exponential operation is executed by a dedicated hardware component, namely the multi-function unit. Ideally, to maximize efficiency, the exponential computation should take place concurrently with the matmul operations handled by the Tensor Cores. To achieve this overlap, we employ synchronization barriers (\texttt{bar.sync} instructions) that enforce an execution order: the GEMM operations (GEMM1 -- $\mathbf{P} \mathbf{V}$ within one iteration, and GEMM0 -- $\mathbf{Q} \mathbf{K}^\top$ from the subsequent iteration) for warpgroup 1 are mandated to complete prior to the GEMM operations of warpgroup 2. This mechanism results in the softmax for warpgroup 1 being processed while warpgroup 2 executes its GEMMs. Afterwards, the two warpgroups alternate their roles, with warpgroup 2 carrying out softmax as warpgroup 1 takes up GEMM operations—this alternating pattern is referred to as ``pingpong'' scheduling. This approach is depicted in~\cref{fig:pingpong_scheduling}. While the real-world execution of pingpong scheduling is not always as orderly as shown in the figure, our empirical results indicate a noticeable performance boost (for example, throughput increases from 570 TFLOPS to 620–640 TFLOPS for FP16 forward passes with a head dimension of 128 and sequence length 8192).

\paragraph{Attention variants} In the cases of multi-query attention~\citep{shazeer2019fast} and grouped query attention~\citep{ainslie2023gqa}, we adopt the method used in \textsc{FlashAttention-2}, modifying tensor indexing so as to prevent redundant storage of $\mathbf{K}$ and $\mathbf{V}$ within HBM.

\subsection{Intra-warpgroup overlapping GEMMs and softmax}

Within a single warpgroup, it is possible to execute certain instructions from the softmax operation concurrently with instructions from the GEMMs. Here, we outline a specific method for achieving this overlap.

In the attention algorithm, operations within the inner loop (main loop) have sequential dependencies that impede parallelization within a single iteration. For example, (local) softmax (lines \ref{code-ws:softmax_start} to \ref{code-ws:softmax_end}) relies on the output $\mathbf{S}_i^{(j)}$ of the first GEMM, while the second GEMM takes its result $\widetilde{\mathbf{P}}_i^{(j)}$ as an operand. Indeed, the wait statements in lines \ref{alg:ws_only_gemm1} and \ref{alg:ws_only_gemm2} of \cref{alg:flash3_wgmma_ws_only} serialize the execution of softmax and GEMMs. However, we can break these dependencies by pipelining across iterations through additional buffers in registers. Pursuing this idea, we propose the following two-stage\footnote{Note that the number of stages of the overlapping scheme is bounded by, but need not equal, the number $s$ of stages in the circular SMEM buffer.} GEMM-softmax pipelining algorithm:

\begin{algorithm}[H]
    \caption{\small\label{alg:flash3_wgmma}\textsc{FlashAttention-3} consumer warpgroup forward pass}
    \begin{algorithmic}[1]
      \REQUIRE  Matrices $\mathbf{Q}_i \in \mathbb{R}^{B_r \times d}$ and $\mathbf{K}, \mathbf{V} \in \mathbb{R}^{N \times d}$ located in HBM, key block size $B_c$ with $T_c = \lceil \frac{N}{B_c} \rceil$.
      \STATE Allocate registers according to the number of consumer warps as required.
      \STATE Set $\mathbf{O}_i = (0) \in \mathbb{R}^{B_r \times d}$ and initialize vectors $\ell_i = (0), m_i = (-\infty) \in \mathbb{R}^{B_r}$ on-chip.
      \STATE Synchronize to ensure that $\mathbf{Q}_i$ and $\mathbf{K}_0$ are available in shared memory.
      \STATE Perform matrix multiplication $\mathbf{S}_{\mathrm{cur}} = \mathbf{Q}_i \mathbf{K}_0^T$ with WGMMA. Commit operation and wait for completion.
      \STATE Free buffer stage 0 for $\mathbf{K}$.
      \STATE Calculate $m_{i}$, $\tilde{\mathbf{P}}_{\mathrm{cur}}$, and $\ell_{i}$ from $\mathbf{S}_{\mathrm{cur}}$, then adjust scaling for $\mathbf{O}_i$.
      \FOR{$1 \le j < T_c - 1$}
        \STATE Synchronize until $\mathbf{K}_j$ is accessible in shared memory.
        \label{code:mainloop_start}
        \STATE Compute next scores via $\mathbf{S}_{\mathrm{next}} = \mathbf{Q}_i \mathbf{K}_{j}^T$ using WGMMA, committing the operation but not waiting for it.
        \label{code:first_wgmma}
        \STATE Await $\mathbf{V}_{j-1}$ being present in shared memory.
        \STATE Update outputs with $\mathbf{O}_{i} = \mathbf{O}_{i} + \tilde{\mathbf{P}}_{\mathrm{cur}} \mathbf{V}_{j-1}$ using WGMMA, commit, but do not wait.
        \label{code:second_wgmma}
        \STATE Wait for the WGMMA execution of $\mathbf{Q}_i \mathbf{K}_{j}^T$ to finish.
        \STATE From $\mathbf{S}_{\mathrm{next}}$, update $m_{i}$, $\tilde{\mathbf{P}}_{\mathrm{next}}$, and $\ell_{i}$.
        \label{code:softmax}
        \STATE Await completion of WGMMA $\tilde{\mathbf{P}}_{\mathrm{cur}} \mathbf{V}_{j-1}$ and rescale $\mathbf{O}_i$ accordingly.
        \STATE Release buffer stages $(j\,\%\,s)$ for $\mathbf{K}$ and $(j-1\,\%\,s)$ for $\mathbf{V}$.
        \STATE Replace $\mathbf{S}_{\mathrm{cur}}$ with $\mathbf{S}_{\mathrm{next}}$.
        \label{code:mainloop_end}
      \ENDFOR \STATE Synchronize for availability of $\mathbf{V}_{T_c - 1}$ in shared memory.
      \STATE Perform the final update
      $\mathbf{O}_{i} = \mathbf{O}_{i} + \tilde{\mathbf{P}}_{\mathrm{last}} \mathbf{V}_{T_c - 1}$ using WGMMA. Commit and wait until completion.

\STATE Finalization: Adjust $\mathbf{O}_{i}$ according to the value of $m_{i}$. Derive $L_{i}$ using both $m_{i}$ and $\ell_{i}$.
Save $\mathbf{O}_{i}$ and $L_{i}$ to HBM, storing them as the $i$-th segment of $\mathbf{O}$ and $L$, respectively.
\end{algorithmic}
\end{algorithm}

\cref{alg:flash3_wgmma} serves as an alternative to the consumer pathway detailed in \cref{alg:flash3_wgmma_ws_only}, thereby forming the complete \textsc{FlashAttention-3} method for FP16 computations. Conceptually, WGMMA is used to refer to asynchronous GEMM operations. During the primary loop (spanning lines \ref{code:mainloop_start} through \ref{code:mainloop_end}), the second WGMMA invocation in iteration $j$ (at line \ref{code:second_wgmma}) is executed concurrently with the softmax operations belonging to iteration $j+1$ (at line \ref{code:softmax}).

Although the pipeline structure depicted above suggests possible improvements in performance, several pragmatic challenges arise:

\paragraph{Compiler reordering} Although the pseudocode assumes an idealized execution sequence, the NVCC compiler often rearranges instructions for optimization purposes. Such reordering can interfere with the intended pipeline between WGMMA and other operations, resulting in possible performance degradation or unintended side effects. Examination of the resulting SASS assembly indicates the compiler produces code with overlapped execution as anticipated (refer to Section~\ref{sec:2-stage-sass}).

\paragraph{Register pressure} Sustaining peak efficiency requires careful management of register usage to prevent register spilling. The two-stage pipeline imposes a higher register requirement to store intermediate computation and maintain pipeline state. In particular, holding an additional $\mathbf{S}_{\mathrm{next}}$ in registers increases register consumption by $B_r \times B_c \times \text{sizeof}(\text{float})$ per threadblock. This increased demand can constrain the use of larger block sizes, which themselves necessitate substantial register allocation. Consequently, optimal configurations should be selected through profiling to balance these competing needs.

\paragraph{3-stage pipelining} Building upon the two-stage design, a three-stage pipeline is proposed to further overlap the second WGMMA and softmax computations. While this may enable improved Tensor Core usage, it necessitates a greater number of registers because of the added pipeline stage, complicating the interplay between tile dimensions and pipeline length. A comprehensive explanation and performance analysis of the three-stage strategy are included in ~\cref{sec:3-stage}.

\subsection{Low-precision with FP8}
\label{sec:algofp8}

\textbf{Efficiency: layout transformations.} Executing the forward pass of \textsc{FlashAttention-3} using FP8 precision introduces unique challenges regarding layout compatibility, which are absent in the FP16 context.

Initially, it is important to recognize that the input tensors $\mathbf{Q}$, $\mathbf{K}$, and $\mathbf{V}$ are commonly arranged with continuity along the head dimension. However, in order to comply with the k-major requirement on FP8 WGMMA during the second GEMM, it is necessary for $\mathbf{V}$—specifically, the SMEM tiles of $\mathbf{V}$—to be contiguous along the sequence length axis. Since the TMA load operation cannot alter which dimension is contiguous, we face two main choices: (1) perform a transpose of $\mathbf{V}$ in GMEM beforehand, or (2) carry out a tile-wise transpose within the kernel after transferring the tiles to SMEM. Option (1) can be realized by either (1a) merging the transpose step into the epilogue of a prior stage such as the rotary embedding, or (1b) utilizing a dedicated pre-processing transpose kernel\footnote{An efficient transpose kernel can approach device bandwidth for speed~\citep{colfax_cutlass_transpose_2024}.} to swap the sequence length and head dimension strides. However, approach (1a) poses challenges for integration into existing library architectures, and approach (1b) imposes excessive overhead in memory-constrained situations like inference.

For FP8 \textsc{FlashAttention-3}, we instead choose the second approach. To perform the in-kernel transpose, we exploit the LDSM (\verb|ldmatrix|) and STSM (\verb|stmatrix|) instructions, which enable warps of threads to collaboratively transfer data between SMEM and RMEM, handling 128-byte segments at a time.\footnote{Although the PTX documentation specifies that LDSM/STSM operate on $8 \times 8$ matrices with 16-bit elements \cite[\S 9.7.13.4.15-16]{ptx}, it is possible to pack pairs of 8-bit elements together to make LDSM/STSM compatible with FP8 computations. Nevertheless, the transpose variants of these instructions cannot separate packed 8-bit values, requiring explicit register shuffling between LDSM and STSM steps in order to achieve a true tile-wise transpose; the exact mechanics are omitted here.} LDSM and STSM not only make efficient use of registers—enabling their use within the producer warpgroup—but also facilitate layout transpositions during memory transfers. Further, beginning from the second iteration, the transpose for the upcoming $\mathbf{V}$ tile can be scheduled to occur concurrently with the execution of the two WGMMAs that process the previous $\mathbf{V}$ and the current $\mathbf{K}$ tile.

Secondly, it is evident that, in contrast to FP16, the arrangement of memory for the FP32 accumulator within an FP8 WGMMA does not align with the memory layout expected for operand A when this operand resides in registers. Portions of these memory arrangements are illustrated in \cref{fig:rmem_accum} and \cref{fig:rmem_operand}, showing the ordering in which register entries are assigned per thread. Utilizing byte permutation instructions, we can reformat the first WGMMA's accumulator to a layout that matches the requirements of the second WGMMA, as well as the configuration of the $\mathbf{V}$ tile generated by the in-kernel transpose. In particular, referring to \cref{fig:rmem_accum}, we reorder the registers as follows:
$$\{ \verb|d0 d1 d4 d5 d2 d3 d6 d7| \},$$
and this register sequence is applied to each 8-byte segment. From the perspective of the logical structure of the $\mathbf{P}$ tile, this operation alters the ordering of its columns—for instance, columns $0189$ are relocated to occupy the first four column positions. To ensure WGMMA computes the appropriate output tile, the in-kernel transpose can be configured to produce a corresponding row arrangement in the $\mathbf{V}$ tile.\footnote{This flexibility in the in-kernel transpose removes the need for shuffle instructions, which would otherwise be necessary to reassign registers across threads, as previously discussed in~\citep{colfax_fp8_flashattention_2024}.}

\textbf{Accuracy: block quantization and incoherent processing.} The FP8 (e4m3) numerical format allocates just 3 bits for representing the mantissa and 4 bits for the exponent, which inherently yields greater rounding error when compared to FP16 or BF16. Additionally, in large-scale models, outlier values are common~\citep{dettmers2208llm, sun2024massive}—these can greatly exceed the majority of values in absolute size, complicating the quantization process. A conventional strategy is to implement per-tensor scaling~\citep{micikevicius2022fp8}, where each tensor (such as $\mathbf{Q}$, $\mathbf{K}$, or $\mathbf{V}$) is assigned a single scalar scale factor. In order to decrease the numerical inaccuracies introduced in FP8-based attention mechanisms, we apply the following two methods:

\begin{enumerate}

\item \textbf{Block quantization}: in this approach, a single scalar value is retained for each block. For each of $\mathbf{Q}$, $\mathbf{K}$, and $\mathbf{V}$, the tensor is partitioned into blocks of dimensions $B_r \times d$ or $B_c \times d$, and each block undergoes quantization independently. This process can be seamlessly integrated with preceding operations such as rotary embedding, without introducing performance overhead (as rotary embedding is constrained by memory bandwidth). Given that the \textsc{FlashAttention-3} algorithm inherently operates over blocks, the scaling for each block of $\mathbf{S}$ naturally accommodates block-level quantization without additional computational expenditure.
\item \textbf{Incoherent processing}: to mitigate the influence of outlier values, a random orthogonal matrix $\mathbf{M}$ is used to transform both $\mathbf{Q}$ and $\mathbf{K}$ prior to FP8 quantization. Due to the orthogonality of $\mathbf{M}$, $\mathbf{M} \mathbf{M}^\top = I$, implying that $(\mathbf{Q} \mathbf{M}) (\mathbf{K} \mathbf{M})^\top$ equates to $\mathbf{Q} \mathbf{K}^\top$. As a result, the transformation does not alter the computed attention outputs, but redistributes outlier elements; each entry in $\mathbf{Q} \mathbf{M}$ or $\mathbf{K} \mathbf{M}$ becomes a randomized linear combination of entries from $\mathbf{Q}$ or $\mathbf{K}$, which serves to diminish quantization errors. In implementation, following \citet{chee2024quip} and~\citet{tseng2024quip}, $\mathbf{M}$ is constructed as the product of random diagonal matrices (with entries $\pm 1$) and a Hadamard matrix. This structure allows for multiplication in $O(d \log d)$ rather than $O(d^2)$ time, and it too can be merged with the rotary embedding operation without incurring additional computational cost.
\end{enumerate}

Experimental results, detailed in \cref{sec:numerical_error}, demonstrate that these two strategies can decrease numerical error by as much as 2.6$\times$.

\section{Empirical Validation}
\label{sec:exp}

We implement \textsc{FlashAttention-3} by leveraging CUTLASS~\citep{Thakkar_CUTLASS_2023} primitives like WGMMA and TMA abstractions, and assess its performance with respect to both efficiency and accuracy.

\begin{itemize}

\item \textbf{Benchmarking attention.} We evaluate the execution time of \textsc{FlashAttention-3} across varying sequence lengths and contrast its performance with the standard PyTorch implementation, \textsc{FlashAttention-2}, a Triton-based version of \textsc{FlashAttention-2} (leveraging H100-tailored instructions), and a vendor-optimized cuDNN implementation of \textsc{FlashAttention-2} for H100 GPUs. Our results show that \textsc{FlashAttention-3} achieves up to a 2.0$\times$ speedup over \textsc{FlashAttention-2} and is 1.5$\times$ quicker than the Triton-based \textsc{FlashAttention-2}. The peak throughput of \textsc{FlashAttention-3} reaches 740 TFLOPs/s, corresponding to 75\% of the H100 GPU's peak theoretical TFLOPs/s.

\item \textbf{Ablation study.} We demonstrate that the introduction of algorithmic advancements, such as warp-specialization and GEMM-softmax pipelining, are responsible for the improved speed observed in \textsc{FlashAttention-3}.

\item \textbf{Accuracy of FP8 attention.} Our evaluation establishes that applying block quantization alongside incoherent processing lowers the numerical error in FP8 \textsc{FlashAttention-3} by a factor of 2.6$\times$.

\subsection{Benchmarking Attention}
\label{subsec:benchmark_attn}

We evaluate the execution time of various attention mechanisms utilizing an H100 80GB SXM5 GPU, examining both configurations with and without causal masking and considering head dimensions of 64 and 128, all with FP16 input precision. The performance data, provided in~\cref{fig:benchmark_attn_fwd} and~\cref{fig:benchmark_attn_bwd}, indicate that \textsc{FlashAttention-3} achieves speedups of approximately 1.5-2.0$\times$ in the forward pass and 1.5-1.75$\times$ in the backward pass relative to \textsc{FlashAttention-2}. In comparison to standard attention baselines, \textsc{FlashAttention-3} offers acceleration ranging from 3$\times$ to 16$\times$. Notably, for sequences of moderate to substantial length (1k tokens or more), \textsc{FlashAttention-3} even outperforms the cuDNN library—an extensively tuned closed-source solution for H100 GPUs.

\paragraph{Benchmark settings:} We vary the sequence length as 512, 1k, ..., 16k, and set batch size so that the total number of tokens is 16k. We set the hidden dimension to 2048, and head dimension to be either 64, 128, or 256 (i.e., 32 heads, 16 heads, or 8 heads). To calculate the FLOPs of the forward pass, we use:

\begin{equation*}
  4 \cdot \text{seqlen}^2 \cdot \text{head dimension} \cdot \text{number of heads}.
\end{equation*}

Causal masking leads us to halve this value, reflecting that nearly 50% of the entries are computed. To estimate the FLOPs required for the backward pass, we take the FLOPs of the forward pass and multiply by 2.5. This ratio comes from having 2 matrix multiplications during the forward computation and 5 matrix multiplications in the backward step, owing to recomputation.

Additionally, we evaluate the forward pass runtime for FP8 using comparable configurations. The performance outcomes for headdim 256 are presented in ~\cref{fig:benchmark_attn_fp8}, while comprehensive results can be found in \cref{sec:benchmark_attn_fp8_full}.

\subsection{Ablation Study: 2-Stage Pipelining Experiments}
\label{sec:2-stage-pipelining-experiments}

We perform an ablation study of both the 2-stage WGMMA-softmax pipeline and warp-specialization within the context of non-causal FP16 \textsc{FlashAttention-3}, holding the parameters fixed at $\{\text{batch}, \text{seqlen}, \text{nheads}, \text{hdim}\} = \{ 4, 8448, 16, 128\}$. The findings reported in~\cref{table:ablation_pipelining} demonstrate that introducing our algorithmic enhancements—including asynchronous execution using warp-specialization and concurrent operations between GEMM and softmax—substantially accelerates performance, increasing throughput from 570 to 661 TFLOPs.

\subsection{Numerical Error Validation}
\label{sec:numerical_error}

Given increased scrutiny regarding the numerical error associated with \textsc{FlashAttention}~\citep{golden2024flash}, we conduct a comparative analysis among \textsc{FlashAttention-2}, \textsc{FlashAttention-3}, and a conventional attention mechanism, referencing an implementation in FP64 as ground truth. In order to emulate extreme features and activations typical in LLMs~\citep{dettmers2208llm, sun2024massive}, we sample the elements of $\mathbf{Q}, \mathbf{K}, \mathbf{V}$ according to the following distribution:

\begin{equation*}
  \mathcal{N}(0, 1) + \mathcal{N}(0, 100) \cdot \mathrm{Bernoulli}(0.001).
\end{equation*}

In this setup, each element follows a normal distribution with a mean of zero and a standard deviation of 1; however, for 0.1\% of the elements, we additionally introduce an independent Gaussian component with a standard deviation of 10. The resulting root mean squared error (RMSE) is reported in~\cref{table:numerical_error}. When using FP16, both \textsc{FlashAttention-2} and \textsc{FlashAttention-3} attain RMSEs that are 1.7$\times$ smaller than those of the conventional implementation, due to storing intermediate results (specifically, the softmax) in FP32. For the baseline FP8 attention, per-tensor scaling is applied, with matrix multiplication accumulation carried out in FP32 and intermediate softmax values held in FP16. Leveraging both block quantization and asynchronous processing, \textsc{FlashAttention-3} in FP8 delivers RMSE that is 2.6$\times$ lower than that of this baseline approach.

\section{Dicussion, Limitations, Conclusion}
\label{sec:discussion}

\textsc{FlashAttention-3} showcases how advancements in programming paradigms and hardware capabilities—including asynchrony and reduced precision—can significantly enhance both the performance and reliability of attention mechanisms. Our approach achieves a 1.5-2.0$\times$ acceleration in attention operations relative to \textsc{FlashAttention-2}, while also decreasing FP8 numerical errors by a factor of 2.6 compared with conventional per-tensor quantization. However, several constraints remain that we aim to tackle in future work: tailoring optimizations for LLM inference, adopting a persistent kernel strategy for the FP8 kernel,\footnote{Within our evaluations, FP16 \textsc{FlashAttention-3} benefits from persistent kernel implementation and load balancing, whereas FP8 \textsc{FlashAttention-3} lacks these features. This difference largely accounts for FP8 \textsc{FlashAttention-3}'s suboptimal performance in cases with short sequence lengths and causal masking relative to the FP8 cuDNN kernels.} and investigating the role of low-precision attention in extensive-scale training. While this study centers on Hopper GPUs, we anticipate that the concepts and optimizations introduced will be applicable to a variety of hardware accelerators. Ultimately, we envision that improvements in the speed and fidelity of attention operations will enable new possibilities for tasks involving long contexts.

\end{document}